% Standalone source-record documentation; no research attempts.
\documentclass[11pt]{article}
\usepackage[margin=25mm]{geometry}
\usepackage{fontspec,unicode-math}
\setmainfont{Latin Modern Roman}
\setmathfont{Latin Modern Math}
\usepackage{amsmath,amssymb,mathtools,xurl,hyperref}
\hypersetup{colorlinks=true,urlcolor=blue,linkcolor=blue}
\setcounter{secnumdepth}{0}
\setlength{\parindent}{0pt}
\setlength{\parskip}{5pt}
\setlength{\emergencystretch}{3em}
\newcommand{\sref}[2]{\hyperlink{source:#1:#2}{[S#2]}}
\newcommand{\cataloguescope}[1]{\paragraph{Source record.}#1}
\newcommand{\floor}[1]{\left\lfloor#1\right\rfloor}
\begin{document}
% UnsolvedMath ID 48; source catalogue code SMA-004
\setcounter{section}{47}
\section{Smale's 4th Problem: Integer Zeros of Polynomials}\label{48}
{\small\sffamily UnsolvedMath ID 48 · Computer Science}
\subsection{Definitions and mathematical statement}

\paragraph{Shared notation.}Standard number systems are $\mathbb N,\mathbb Z,\mathbb Q,\mathbb R,\mathbb C$. Specialized notation and index ranges follow the cited source and the definitions below.


For $p\in\mathbb Z[x_1,\ldots,x_r]$, an integer zero is a vector $z\in\mathbb Z^r$ satisfying $p(z)=0$. A decision algorithm must halt and correctly report existence or nonexistence for every permitted input. Its complexity depends on the number of variables, the coefficient encoding and degree representation. In the actual Smale Problem 4, a univariate polynomial is represented by a straight-line program starting from $1$ and $t$ and using addition, subtraction and multiplication; $\tau(f)$ is the smallest number of operations.
\paragraph{Record scope and qualifications.}
The catalogue title is retained as metadata. Named certificates, contracts, bounds or frameworks in that title are not asserted as results of the original source extract.
The supplied title and decision formulation do not agree with Smale Problem 4. The original paper asks for a polynomial bound on the number of distinct integer roots in terms of straight-line complexity, rather than the decision task quoted below. Its Problem 5 discussion also records undecidability when the number of variables is unrestricted. The catalogue formulation is preserved, without substituting the root-count conjecture or claiming a universal decision algorithm.
\paragraph{Statement recorded in the supplied source.}
Find efficient algorithms for deciding whether a polynomial with integer coefficients has an integer root. \sref{48}{2}
\paragraph{Source status.}
The supplied catalogue labels this source record open; that metadata was checked against the supplied snapshot on 2026-10-11. The label is preserved for provenance and is not an independent verification of present open status or mathematical correctness.
\subsection{Short English statement}
The catalogue asks for efficient integer-root decisions. Its title points to a different Smale question about how many integer roots a univariate polynomial can have.
\subsection{Sources}
\begin{itemize}
\item[\hypertarget{source:48:1}{S1}] Ulam AI and the UnsolvedMath Contributors, UnsolvedMath: A Curated Collection of Open Mathematics Problems, record 48 (SMA-004). Catalogue curation and metadata: CC BY 4.0. The source extract is retained; definitions, English summaries and qualifications are editorial additions.
\url{https://huggingface.co/datasets/ulamai/UnsolvedMath}.
\item[\hypertarget{source:48:2}{S2}] Steve Smale, Mathematical Problems for the Next Century, second version (7 August 1998), Problems 4 and 5, pp. 5--6. Consulted for the attribution and scope comparison; the catalogue decision wording differs.
\url{https://www.fim.uni-passau.de/fileadmin/dokumente/fakultaeten/fim/lehrstuhl/muller/SmaleProblems1998.pdf}.
\item[\hypertarget{source:48:3}{S3}] Steve Smale, Mathematical Problems for the Next Century, Mathematical Intelligencer 20(2), 7--15 (1998).
\url{https://doi.org/10.1007/BF03025291}.
\end{itemize}
\paragraph{Preparation and review.}
The original source and UnsolvedMath contributors retain attribution for the source extract. Codex (GPT-6; exact serving revision not exposed) assisted with definitions, English summaries, scope notes and formatting. No new human mathematical authorship or independently verified solution is asserted. The supplied catalogue credits: Stephen Smale. Permanent source ID 48 and code SMA-004 checked against the complete supplied registry and current repository. Distinct upstream records, including repeated questions and broad or damaged extracts, are retained for complete coverage of the source CS category. A repeated record is not a new independent conjecture. Complete mathematical-equivalence and current-literature review has not been performed. Preparation and final review follow contribution prompts C and D, with source gaps disclosed. Structural and rendering checks do not establish a new theorem or open status.
\label{end:48}
\end{document}
